\section*{Chapter \ref{ch:in:medium-frequency-proprietary-shops} --- Medium-Frequency Proprietary Shops}

\begin{solution}{exo:in:medium-frequency-proprietary-shops:1}
$65+84=149$ of 3\,249 firms, 4.6\%, employing 81.7\% of the representatives.
\end{solution}

\begin{solution}{exo:in:medium-frequency-proprietary-shops:2}
$0.4\times(500\,000-150\,000)=\$140\,000$; on \$150\,000 the costs take everything and the trader earns nothing.
\end{solution}

\begin{solution}{exo:in:medium-frequency-proprietary-shops:3}
$(1\,000.5/10.5)^{-0.6}=0.065$: 6.5\%.
\end{solution}

\begin{solution}{exo:in:medium-frequency-proprietary-shops:4}
Exit 5.6\% and entry 3.2\% a year on average: net shrinkage of about 2.4\% a year.
\end{solution}

\begin{solution}{exo:in:medium-frequency-proprietary-shops:5}
With $\alpha<1$ the tail's mean is infinite: the expected size of the largest firm grows faster than the population,
so the largest firms' share of the total does not shrink as firms are added.
\end{solution}

\begin{solution}{exo:in:medium-frequency-proprietary-shops:6}
From about 300 representatives up, and most at the top (6.5\% predicted above 1\,000 against 4.9\% observed): the very
largest firms are fewer than a pure power law implies, as if their size had a ceiling.
\end{solution}

\begin{solution}{exo:in:medium-frequency-proprietary-shops:7}
2020: $\alpha=0.616$ (standard error 0.015) against 0.601 in 2024. The difference, 0.015, is below two standard errors of
the difference, $2\sqrt{0.015^2+0.015^2}=0.042$: no change the data can detect.
\end{solution}

\begin{solution}{exo:in:medium-frequency-proprietary-shops:8}
The unit is registered representatives, not staff, and the count is of firms, not jobs: most registered people work at
the 4.6\% of firms with 500 or more. A typical job is at a large firm, even though a typical firm is small.
\end{solution}

\begin{solution}{pb:in:medium-frequency-proprietary-shops:1}
\begin{enumerate}
\item Systematic trading with positions of minutes to days; slower than high-frequency trading, faster than most
systematic funds.
\item Its costs are impact and fees, not being first; it needs good data and execution but not tick-to-trade engineering.
\item Registered representatives; a firm without customers registers few of its staff.
\item 103, of which 93 small (150 representatives or fewer) and 10 mid-size; none large.
\item 46.8\% and 78.7\%.
\item $p_i=((\mathrm{lo}_i-\tfrac12)/(x_{\min}-\tfrac12))^{-\alpha}-((\mathrm{hi}_i+\tfrac12)/(x_{\min}-\tfrac12))^{-\alpha}$.
\item $0.60$, standard error 0.015.
\item Lognormal 5.1, Pareto 21.2: the lognormal fits the bins better.
\item It rises from 0.60 at 11 to 0.65 at 51.
\item Heavy concentration: 9.4\% of representatives at small firms, 81.7\% at the largest 4.6\%.
\item 5.6\% and 3.2\%; 185 left and 136 entered in 2024.
\item Several entities or persons trading under one licence; all investment-firm requirements met, and every
participant's rights and obligations allocated and documented.
\item Customers pay fees for evaluations and are promised profit shares; its revenue is their fees, and they are its
customers, not its staff.
\item The CFTC alleged in 2023 that over 135\,000 customers paid at least \$310 million in fees; in May 2025 a special master
recommended dismissal with prejudice on a motion for sanctions against the CFTC, and the court dismissed the case with
prejudice and awarded fees. The allegations were not proved.
\item \$260\,000, \$20\,000 and nothing.
\item $\alpha$ about 0.60 (0.015), between 0.60 and 0.65 as the lower bound moves; 9.4\% of representatives at small firms.
\item The capital is the partners'; pay follows P\&L; controls are lighter; the firm may not last.
\item Whose capital and how large; the split, the costs charged and loss carry-forward; who checks risk and trades.
\item About 5.6\%, one in eighteen.
\item The offer is a share in a small business more than a salary: its value depends on the split, the costs and the firm's
survival. Weigh it for what it teaches and against a salaried alternative's lower variance.
\end{enumerate}
\end{solution}

\begin{solution}{iq:in:medium-frequency-proprietary-shops:1}
The edge moves from speed to prediction: signals on minutes to days, with costs dominated by impact and fees; research
leans on statistics and data, infrastructure on reliable data and execution rather than microsecond latency; capacity
rises.
\iqlookfor{horizon drives edge, cost and infrastructure.}
\end{solution}

\begin{solution}{iq:in:medium-frequency-proprietary-shops:2}
$0.5x=150\,000+0.1x$, so $x=\$375\,000$ of net P\&L. Above it the split pays more.
\iqlookfor{a break-even and the variance each carries.}
\end{solution}

\begin{solution}{iq:in:medium-frequency-proprietary-shops:3}
Maximise the binned likelihood of a Pareto tail above a chosen lower bound; get a standard error from the observed
information; compare with alternatives (lognormal) by likelihood or chi-square, and check stability as the lower bound
moves.
\iqlookfor{binned maximum likelihood, alternatives and the choice of $x_{\min}$.}
\end{solution}

\begin{solution}{iq:in:medium-frequency-proprietary-shops:4}
All of it: every participant's losses, capital requirements and conduct are the licence holder's; so are operational
failures and market-abuse risk. It needs pre-trade limits per participant, real-time monitoring and clear documentation.
\iqlookfor{the licence holder carries the risk.}
\end{solution}

\begin{solution}{iq:in:medium-frequency-proprietary-shops:5}
$(5-2)=3$ basis points a day on 100\% turnover: about $3\times252=756$ basis points, 7.6\% a year before financing. Impact
grows with size, so the 5 basis points shrink as capital rises; capacity ends where marginal edge equals marginal cost.
\iqlookfor{edge net of cost times turnover, and capacity.}
\end{solution}

\begin{solution}{iq:in:medium-frequency-proprietary-shops:6}
Compare fee revenue with profit-share payouts and with trading P\&L on real accounts; ask what share of customers pass,
whether funded accounts trade in a market or against the operator, and how payouts are funded. A business whose payouts
are small against fees is selling evaluations.
\iqlookfor{follow the revenue.}
\end{solution}
